\chapter{Coupled Simulation}
Now we couple locally mass conserved MFEM and FVM method together in this chapter.
We discuss implementation details and coupling strategy combining MFEM, FVM and eutectic phase
package.
The remaining uncertainty of porosity $\phi$ will be fulfilled by phase package
evaluation. Before we throw ourselves into the fully coupled problem.
We test another porosity evolution scenerio with a more straightforward prescribed melting.
Then we will show an example combining MFEM, FVM and eutectic phase package
alltogether and perform coupled computation.

\section{Coupled Algorithm}
Let $U^n=(C^n, H^n)$ be the solution to the transport system at time $t^n$. 
Let $V^n=(\check{\mathbf{v}}_{h,r}^n,\check{\mathbf{v}}_{h,s}^n,\check{q}_{h,l}^n,\check{q}_h^n)$ 
denote the solution to the Darcy-Stokes system at the same time. Moreover, let $W^n=(\phi^n, \check{T}^n, c_{s}^n, c_{l}^n, \phi_1^n)$ 
denote the output of the phase calculation at $t^n$. We formulate an overall solution methods.
\begin{algorithm}[htb]
\caption{Sequential Coupled Solver\label{alg:coupled}}
%\SetAlgoLined
 Let $U^0$ be given as initial data.

\begin{algorithmic} [1]
 \For{\rm Time level $n=0,1,\ldots,n_{\max}-1$}
		  \vspace{0.3cm}
\State Given $U^n$, compute $W^n$ by the eutectic phase package.
		  \vspace{0.2cm}
\State Given $\phi_f^n \in W^n$, solve the Darcy-Stokes system 
		  \eqref{eq:dimensionlessdarcy_flow_sc}--\eqref{eq:dimensionlesscompaction_flow_sc} for $V^{n+1}$.
		  \vspace{0.2cm}
		  \State Given $W^n$ and $V^{n+1}$, compute effective 
		  velocity~\eqref{eq:definitioneffvel} and phase averaged velocity as
\begin{equation*}
\bold{v}_{\text{eff}} = \frac{\phi_l \bold{v}_l + D_i\phi_s \bold{v}_s}{\phi_l + \phi_s D_i}, \quad \overline{\bold{v}} = \phi \mathbf{v}_r + \mathbf{v}_s.
\end{equation*}
		  \vspace{0.2cm}
        \State Time stepping the transport system~\eqref{eq:transportFV} to next time $n+1$ for $U^{n+1}$.
%     \State Given $U^{n+1}$, compute $W^{n+1}$ by the eutecetic phase package.
%     \State Given $\phi_f^{n+1}$ in $W^{n+1}$, solve the Darcy-Stokes system \eqref{eq:darcy_flow_sc_fem}--\eqref{eq:compaction_flow_sc_fem} for $V^{n+1}$.
		  \vspace{0.3cm}
\EndFor
\end{algorithmic}
\end{algorithm}
The coupled solver is initialized from the initial condition on the transport
variables $U^0$ specifying the total initial mass and energy. The eutectic 
phase package gives $W^0$. Once we know $W^0$, we have phase related variables
porosity and individual concentrations.

In Algorithm~\ref{alg:coupled}, $U^{n+1}$ is computed using fixed in time values for the effective velocities $\mathbf{v}_{\text{eff}}$ and phase 
averaged velocity $\overline{\mathbf{v}}$, 
 pre-calculated from the Darcy-Stokes part, and the partition functions $\kappa$ and $\lambda$. 
 These quantities are only updated before commencing the next time step.

\subsection{Time stepping}
In this section, we briefly talk 
about suitable time integration technique. 
We will focus on explicit scheme for two reasons
\begin{itemize}
		  \item Phase package cannot be
evaluated with long time intervals.
\item Jacobian is hard to evaluate since a MFEM solver is involved.
\end{itemize}
While forward Eular is sufficient in stability analysis,
it only provides first order accuracy. Therefore, we use strong-stability 
preservation (SSP) scheme proposed by~\cite{doi:10.1137/S003614450036757X},
which is a higher order time stepping scheme while preserving 
the stability properties of forward Euler.
We start with first order forward Euler as
\begin{equation}
u^{n+1} = u^n + \Delta t \tilde{\mathbf{F}}(u^n),
\end{equation}
where $\tilde{\mathbf{F}}(u^n)$ represents a spatial discretization.

We state a multi-stage SSP Runge-Kutta  method as
\begin{equation}\label{eq:multistageSSP}
\begin{aligned}
		  u^{0} &= u^n ,\\
		  u^{(i)}   &= \sum_{j=0}^{i-1} \alpha_{i,j}u^{(j)} + \Delta t \beta_{i,j} \tilde{\mathbf{F}}(u^{(j)}), \quad 1\leq i \leq m\\
		  u^{n+1} &= u^m,
\end{aligned}
\end{equation}
where $\sum_{j}^{i-1} \alpha_{i,j} = 1$ for consistency. 
In our computation, we have second order accuracy in MFEM solver, and a MLWENO(3,2)
reconstruction in FVM. Hence the highest order of time stepping accurcay is three.
We write second and third order SSP Runge-Kutta explicit as
\begin{equation}
\begin{aligned}
		  u^{0}   &= u^n\\
		  u^{(1)} &= u^{(0)} + \Delta t \tilde{\mathbf{F}}(u^{(0)})\\
		  u^{n+1} &= u^{(0)} + \frac{1}{2} \Delta t \tilde{\mathbf{F}}(u^{(0)}) + \frac{1}{2} \Delta t \tilde{\mathbf{F}}(u^{(1)}),
\end{aligned}
\end{equation}
and 
\begin{equation}
\begin{aligned}
		  u^{0}   &= u^n ,\\
		  u^{(1)} &= u^{(0)} + \Delta t \tilde{\mathbf{F}}(u^{(0)}),\\
		  u^{(2)} &= \frac{3}{4} u^n + \frac{1}{4} u^{(1)} + \frac{1}{4} \Delta t \tilde{\mathbf{F}}(u^{(1)})\\
		  u^{n+1} &= \frac{1}{3} u^{(0)}+\frac{2}{3}u^{(2)} + \frac{2}{3}\Delta t \tilde{\mathbf{F}}(u^{(2)}).
\end{aligned}
\end{equation}

\subsection{Boundary Conditions}
In practical implementation, we prescribe boundary conditions to each dofs on the boundary.
Each dofs will be associated with two types of boundary conditions. 

\section{Prescribed Melting Problem}
Before fully coupled computation, 
we illustrate a simpler model coupling case coupling MFEM and FVM
solver. Instead of deciding porosity with phase package computation,
we prescribe the melting rate or the change of porosity as a 
fixed function distributed in space as
\begin{equation}\label{eq:gamma}
\phi(t) = \Gamma(\mathbf{x}).
\end{equation}
We write conservation of liquid mass with the prescribed 
melting term $\Gamma$~\eqref{eq:gamma} as
\begin{equation}
		  \frac{\partial \phi}{\partial t} + \nabla \cdot (\phi \mathbf{v}_l) = \Gamma,
\end{equation}
where diffusion effect is being dropped for simplicity.

We perform a quasi-1D column test on a $2\times40$ rectangular physical domain $\Omega \in [-0.05,0.05] \times [-0.4,0]$.
We prescribe the following boundary conditions for MFEM and FVM solver.
\begin{itemize}
		  \item Bottom side;
					 \begin{itemize}
								\item MFEM. We prescirbe solid velocity at the bottom of the 
physical domain, indicating supplying of mantle. 
\begin{equation*}
\check{\mathbf{v}}_s = \mathbf{v}_{\text{upwelling}}.
\end{equation*}
We prescribe zero liquid velocity
indicating no melting penetrating bottom of the computational domain;
\item FVM. We prescribe a fixed inflow boundary condition, supplying composition 
		  and enthalpy into system.
					 \end{itemize}
		  \item Left and right side;
					 \begin{itemize}
								\item MFEM. We prescribe natural boundary conditon giving zero
										  traction to the Stokes equation in tangential direction.
										  We prescribe zero flux for dofs in normal direction for 
										  the Stokes equation and Darcy equation.
										  \item FVM. We prescribe noflow (zero flux) condition.
					 \end{itemize}
					 \item Top side;
					 \begin{itemize}
                \item MFEM. We prescribe natural boundary conditoin given zero traction
					 and zero pressure to Stokes and Darcy part respectively.
                \item FVM. We set free outflow boundary if it is recognized as outflow.
					 We set zero if the velocity reversed to inflow boundary condition.
					 \end{itemize}
\end{itemize}
For simplicity and efficiency in computation, we use $1\times 3$ and $1\times 2$ for our
ML-WENO(3,2) stencils, since we don't need resolution and high order accuracy in x direction.
We show two different melting source scenarios. One of which is a piece wise constant
stair function:
\begin{equation}
\Gamma_1(z) = \left \{ 
\begin{aligned}
&1e-3, \hspace{0.2cm} -0.2 < z\leq 0,\\
&5e-4,\hspace{0.2cm} -0.3 < z\leq -0.2,\\
&0, \hspace{1.13cm}   -0.4 \leq z\leq -0.3,
\end{aligned}
		  \right.
\end{equation}
and one with melting source trapped in the middle of the melting column as
\begin{equation}
\Gamma_2(z) = \left \{ 
\begin{aligned}
&0, \hspace{1.13cm} -0.2 < z\leq 0,\\
&5e-4,\hspace{0.2cm} -0.3 < z\leq -0.2,\\
&0, \hspace{1.13cm}   -0.4 \leq z\leq -0.3,
\end{aligned}
		  \right.
\end{equation}
We assign different initial conditions to these two scenarios as
\begin{equation}
\phi_1(z,0) =  0.0 , \quad z \in [-0.4,0.0],
\end{equation}
and
\begin{equation}
\phi_2(z,0) = \left \{ 
\begin{aligned}
&1e-2, \hspace{0.2cm} -0.2 < z\leq 0,\\
&0, \hspace{1.13cm}   -0.4 \leq z\leq -0.2.
\end{aligned}
		  \right.
\end{equation}



\subsection{Fully coupled Computation}
